El instrumento de validación incluyó dos preguntas abiertas (PA1 y PA2) orientadas a recoger observaciones cualitativas de los expertos sobre los elementos faltantes o poco desarrollados del modelo y sus sugerencias de mejora. A partir de un análisis temático de las respuestas, se identificaron seis categorías emergentes que sintetizan los principales señalamientos. La Tabla~\ref{tab:categorias-cualitativas} presenta estas categorías junto con los expertos que las mencionaron.

\begingroup
\scriptsize
\setlength{\LTleft}{0pt}
\setlength{\LTright}{0pt}
\setlength{\tabcolsep}{4pt}
\renewcommand{\arraystretch}{1.4}

\begin{longtable}[H]{|
   >{\Centering\hspace{0pt}}p{0.8cm}|
   >{\RaggedRight\hspace{0pt}}p{3.8cm}|
   >{\RaggedRight\hspace{0pt}}p{7.8cm}|
   >{\Centering\hspace{0pt}}p{2.3cm}|
   }

   \caption{Categorías emergentes del análisis cualitativo de preguntas abiertas}
   \label{tab:categorias-cualitativas}                                                                                                                                                                   \\
   \hline
   \textbf{Id.}         &
   \textbf{Categoría}   &
   \textbf{Descripción} &
   \textbf{Expertos}                                                                                                                                                                                     \\
   \hline
   \endfirsthead

   \hline
   \multicolumn{4}{|c|}{{\bfseries \tablename\ \thetable{} -- Continuación}}                                                                                                                             \\
   \hline
   \textbf{Id.}         &
   \textbf{Categoría}   &
   \textbf{Descripción} &
   \textbf{Expertos}                                                                                                                                                                                     \\
   \hline
   \endhead

   \hline
   \multicolumn{4}{|r|}{\textit{Continúa en la siguiente página...}}                                                                                                                                     \\
   \hline
   \endfoot

   \hline
   \multicolumn{4}{|p{14.5cm}|}{\footnotesize\textit{Acrónimos utilizados: Identificador (Id.), Categoría (C) y Expertos (E).}}                                                                          \\
   \hline
   \endlastfoot

   C1                   & Validación empírica en contextos reales      & Necesidad de aplicar el modelo en organizaciones reales para obtener evidencia de su efectividad                   & E1, E5, E6 \\
   \hline
   C2                   & Preparación y formación del Social Guide     & Ausencia de rutas de formación para que una persona nueva pueda asumir el rol de Social Guide                      & E3, E4     \\
   \hline
   C3                   & Protocolos de resolución de conflictos       & Falta de mecanismos explícitos de mediación o escalamiento cuando surgen desacuerdos entre roles                   & E3         \\
   \hline
   C4                   & Gestión de interrupciones y estado de flujo  & Riesgo de que las intervenciones sociales interrumpan períodos de trabajo profundo y afecten la productividad      & E4         \\
   \hline
   C5                   & Estrategias tácticas y catálogo de dinámicas & Carencia de un repositorio de dinámicas concretas para mitigar community smells y de métricas de impacto asociadas & E4, E5     \\
   \hline
   C6                   & Adopción gradual y marco de madurez          & Necesidad de niveles progresivos de implementación para contextos organizacionales con distintos grados de madurez & E2, E5, E6 \\
\end{longtable}
\endgroup

\subsubsection{Análisis de PA1: Elementos faltantes o poco desarrollados}

La primera pregunta abierta solicitó a los expertos identificar elementos del modelo que consideraran ausentes o insuficientemente desarrollados. Del análisis de las respuestas emergieron cuatro ejes temáticos principales.

\textbf{Validación empírica (C1).} La categoría con mayor recurrencia fue la necesidad de validación empírica en contextos reales, señalada por tres de los seis expertos. E1 indicó que \textit{``el modelo necesita ponerse en práctica en contextos reales [\ldots] para obtener pruebas empíricas que permitan evaluar su efectividad''}. De forma convergente, E5 señaló la ausencia de \textit{``validación empírica en contextos reales''} como una de las principales carencias, mientras que E6 reconoció que \textit{``el modelo no ha sido validado empíricamente en múltiples organizaciones, lo que significa que aún no existe evidencia sólida de su efectividad en distintos contextos reales''}. Este hallazgo es consistente con la valoración cuantitativa de P11 ($\mu = 4{,}00$; $\sigma = 0{,}82$), que reflejó la mayor dispersión del instrumento precisamente en el aspecto de factibilidad de implementación.

\textbf{Preparación del Social Guide (C2).} E3 propuso que, dado que el Social Guide constituye el rol novedoso del modelo, sería conveniente \textit{``explicar también cómo una persona nueva dentro de la organización podría prepararse para cumplir con este rol''}. Este señalamiento complementa la dispersión observada en P10 ($\mu = 4{,}50$; $\sigma = 0{,}76$), que indicaba la necesidad de afinar la definición del perfil. Por su parte, E4 señaló la carencia de \textit{``protocolos de ejecución para las intervenciones sociales que protejan la productividad individual''}, lo cual también se vincula con la operacionalización del rol.

\textbf{Gestión de interrupciones y estado de flujo (C4).} E4 advirtió que \textit{``una interrupción en un momento crítico puede duplicar el tiempo de recuperación del enfoque''} y planteó la necesidad de definir \textit{``si las actividades del Social Guide deben ser asíncronas o en bloques de tiempo que no fracturen la jornada de desarrollo''}. Esta observación introduce una dimensión operativa que el modelo actual no aborda de forma explícita.

\textbf{Adopción gradual (C6).} E6 observó que el modelo \textit{``solo puede aplicarse si existen prerrequisitos como cultura mínimamente saludable y compromiso ejecutivo''}, y señaló que \textit{``faltarían estrategias de adopción gradual en entornos inmaduros y versiones `light' del modelo para contextos menos favorables''}. E2, si bien valoró positivamente los elementos del modelo, reconoció la necesidad de \textit{``una estandarización oficial del modelo''} dado que \textit{``las empresas lo apliquen de forma diferente acorde a sus necesidades''}.

\subsubsection{Análisis de PA2: Sugerencias de mejora}

La segunda pregunta abierta solicitó a los expertos formular sugerencias de mejora con fundamento teórico o práctico. Las propuestas se agrupan en tres ejes temáticos.

\textbf{Estudios piloto y métricas complementarias (C1, C5).} E1 recomendó \textit{``aplicarlo en equipos y organizaciones para obtener pruebas empíricas que permitan evaluar su efectividad y, a partir de esa experiencia, identificar oportunidades de mejora''}. E5 amplió esta sugerencia proponiendo \textit{``validación empírica mediante estudios piloto en distintas organizaciones''} e incorporar \textit{``indicadores más objetivos (por ejemplo, análisis longitudinal y métricas de rotación o ausentismo) para complementar el Health Score''}. E4 planteó la necesidad de un \textit{``repositorio de dinámicas validadas''} que especifique \textit{``cómo cada una impacta directamente en métricas de rendimiento técnico''}.

\textbf{Protocolos de mediación y gestión del tiempo (C3, C4).} E3 sugirió \textit{``definir protocolos de mediación o escalamiento para estas situaciones cuando el Social Guide está involucrado''}. E4 complementó esta línea proponiendo regular el \textit{timing} de las intervenciones y considerar el impacto de \textit{``las actividades sociales síncronas durante las `Golden Hours' de desarrollo''}, advirtiendo que \textit{``una interrupción social en un momento de alta concentración puede causar una pérdida de hasta 20--30 minutos de re-enfoque''}.

\textbf{Marco de madurez y escalabilidad (C6).} E6 propuso \textit{``definir niveles progresivos de implementación (Inicial, Intermedio, Avanzado)''}, argumentando que \textit{``los marcos de cambio organizacional (Kotter, CMMI, SAFe) muestran que los modelos escalables requieren niveles graduales para contextos heterogéneos''}. E5 sugirió, en línea con lo anterior, \textit{``desarrollar un modelo de escalabilidad para múltiples equipos, apoyado en marcos como Scrum@Scale''}.

En conjunto, el análisis cualitativo refuerza y profundiza los hallazgos del análisis cuantitativo: las áreas con mayor dispersión estadística (P4, P10 y P11) coinciden con las categorías temáticas más recurrentes en las respuestas abiertas, lo que otorga mayor robustez a la identificación de las áreas de mejora del modelo.
